\chapter{Personal Financial Statements and Ratio Analysis}

Before deciding where to go, I need to know where I am. This chapter presents my personal balance sheet, income summary, and four financial ratios benchmarked against standard guidelines.

\section{Personal Balance Sheet (as of June 2026)}

\begin{table}[h!]
\caption{Personal Balance Sheet (as of June 2026)}
\label{tab:balance_sheet}
\centering
\begin{tabular}{@{}lrlr@{}}
\toprule
\textbf{Assets} & \textbf{Value (VND)} & \textbf{Liabilities \& Net Worth} & \textbf{Value (VND)} \\
\midrule
\textit{Liquid Assets} & & \textit{Liabilities} & \\
Cash on Hand (ABBank) & 8,000,000 & Current Debt & 0 \\
\textit{Investments} & & Long-term Debt & 0 \\
Stock Portfolio (HOSE) & 28,000,000 & \textbf{Total Liabilities} & \textbf{0} \\
Bitcoin (Crypto)       & 20,000,000 & & \\
\textit{Personal Use Assets} & & & \\
Xiaomi 15 Smartphone   & 11,000,000 & & \\
Datbike Quantum S2     & 43,000,000 & \textbf{Net Worth} & \textbf{126,500,000} \\
ROG Zephyrus G14 Laptop & 16,500,000 & & \\
\midrule
\textbf{Total Assets} & \textbf{126,500,000} & \textbf{Total Liabilities \& NW} & \textbf{126,500,000} \\
\bottomrule
\end{tabular}
\end{table}

\textit{Note.} All personal use assets are valued at current resale market prices.

Total assets are 126,500,000~VND with zero debt. The largest single asset is the Datbike Quantum S2 electric motorcycle, currently valued at 43,000,000 VND. Of the 56,000,000 VND in ``liquid'' assets, only 8,000,000 VND is cash; the remaining 48,000,000 VND is in stocks and Bitcoin, which can drop in value. This is the key weakness addressed in the ratio analysis below.

\section{Income and Expense Summary}

Monthly gross income is 9,500,000~VND from three sources: a fixed internship allowance of 3,000,000 VND from kNowHow Consultant, a monthly family allowance of 4,000,000 VND, and average freelance income of 2,500,000 VND. Monthly out-of-pocket expenses are 5,189,865~VND, leaving a monthly surplus of 4,310,135~VND.

\section{Financial Ratio Analysis}

\begin{enumerate}
  \item \textbf{Savings Ratio:}
    \begin{equation*}
      \text{Savings Ratio} = \frac{4{,}310{,}135}{9{,}500{,}000} \approx 45.4\%
    \end{equation*}
    Benchmark: $>$20\%. Result: 45.4\%---strong, but dependent on not paying rent. If rent were included (4,000,000 to 7,000,000~VND/month), this ratio would be near zero.

  \item \textbf{Liquidity Ratio:}
    \begin{equation*}
      \text{Liquidity} = \frac{56{,}000{,}000}{5{,}189{,}865} \approx 10.8 \text{ months}
      \quad \Rightarrow \quad
      \text{Cash-only} = \frac{8{,}000{,}000}{5{,}189{,}865} \approx 1.5 \text{ months}
    \end{equation*}
    Benchmark: 3--6 months. The headline figure (10.8) looks fine, but cash-only (1.5 months) is below the safe minimum. Building cash reserves to at least 15,000,000~VND is the most urgent action in the near-term plan.

    The distinction between the two numbers matters in practice. If markets fall sharply and I need emergency cash, I cannot rely on selling stocks or Bitcoin at their current stated values---both can lose 30 to 50\% in a downturn, at exactly the moment when cash is most needed. A real emergency fund must be in cash. Until the 15,000,000~VND floor is reached, no additional capital should go into stocks or crypto.

  \item \textbf{Solvency Ratio:}
    \begin{equation*}
      \text{Solvency} = \frac{126{,}500{,}000}{126{,}500{,}000} = 100\%
    \end{equation*}
    Every asset is fully owned---zero debt. Ideal starting point for early wealth-building.

  \item \textbf{Debt Service Ratio:}
    \begin{equation*}
      \text{DSR} = \frac{0}{9{,}500{,}000} = 0\%
    \end{equation*}\nopagebreak[4]
    No monthly debt payments. The full 4,310,135~VND surplus goes to saving and investing. Future debt (vehicle 2033, mortgage 2040) is analyzed in Phases~6 and~7.
\end{enumerate}
